\section{The gap is asymptotically deterministic}\label{sec:variance}

This section shows that the gap $\Dq=\LPq-\Tq$ has variance $O(q^{-1/6+\eps})$ for every
$\eps>0$. The proof has three parts. First, with high probability every optimal
assignment uses only pairs of deficit $O(q^{-1/2+\eps})$, and every optimal solution of the
fractional relaxation is supported on teams of deficit $O(q^{-1/2+\eps})$
(Propositions~\ref{prop:variance-assignment} and~\ref{prop:variance-lpsupport}); both rest
on an expansion property of the market that holds uniformly over all sets of positions
(Lemmas~\ref{lem:variance-expansion} and~\ref{lem:variance-growth}). Second, when the pair
deficits are written as functions of independent exponential variables, the gradient of
$\Tq$ is the vector of deficits of the pairs used by the optimal assignment, and the
gradient of $\LPq$ is the vector of deficits weighted by the pair loads
(Lemma~\ref{lem:variance-gradient}). Third, the Poincar\'e inequality for the exponential
measure (Lemma~\ref{lem:variance-poincare}) bounds the variance by the expected squared
gradient, which is at most the largest deficit in use times the total deficit; the total
deficit is $O(q^{1/3})$ in expectation by a theorem of Frieze, Pegden and Tkocz
(Proposition~\ref{prop:variance-fpt}).

Throughout the section logarithms are natural, and a team $e=(a,b,c)$ always has $a\in A$,
$b\in B$, $c\in C$. A \emph{pair} is two agents from different sectors. The value of a team
is $v(e)=3-\Defi(e)$, where $\Defi(e)=\defi(a,b)+\defi(b,c)+\defi(a,c)$. For a set
$\mathcal M$ of disjoint teams we write $\Defi(\mathcal M)=\sum_{e\in\mathcal M}\Defi(e)$
and call it the deficit of $\mathcal M$; it is the sum of the deficits of the pairs
\emph{used} by $\mathcal M$, a pair being used if some team of $\mathcal M$ contains both
of its agents. A perfect assignment is determined by the set of $3q$ pairs it uses. An
\emph{optimal assignment} is a perfect assignment of least deficit, so that
$\Phiq=3q-\Tq$ is the deficit of an optimal assignment. When the size of the market has to
be shown we write $T_n$ and $\Phi_n$ for the market with $n$ agents in each sector.

\subsection{Statement}\label{sec:variance-statement}

The only input that is not elementary is the following.

\begin{proposition}[Frieze, Pegden and Tkocz {\cite[Theorem~2(b)]{FPT20}}]\label{prop:variance-fpt}
Consider the market with $n$ agents in each sector, the cost of a team being its deficit,
that is, the sum of three pair costs, where the $3n^2$ pair costs are independent and
uniform on $[0,1]$. There is a constant $c_2$ such that the greedy algorithm of
\cite{FPT20}, run until $n-n^{1/4}$ disjoint teams have been chosen, finds a solution of
expected cost at most $c_2n^{1/3}$; the last step of the algorithm, which completes the
teams chosen greedily to a perfect assignment, may be carried out by an arbitrary
completion. Consequently
\begin{equation}\label{eq:variance-C0}
C_0:=\sup_{n\ge1}\frac{\E\Phi_n}{n^{1/3}}<\infty .
\end{equation}
\end{proposition}

\begin{proof}[Proof of \eqref{eq:variance-C0} from the cited statement]
The completed output of the greedy algorithm is a perfect assignment, so $\Phi_n$ is at
most its deficit. Whether or not the bound $c_2n^{1/3}$ is read as including the completion,
the completion adds at most $n^{1/4}+1$ teams of deficit at most $3$ each, so
\[
\E\Phi_n\le c_2n^{1/3}+3n^{1/4}+3\le(c_2+6)\,n^{1/3}.
\]
If the cited statement is read as holding only for large $n$, then $C_0$ is still finite,
because $\Phi_n\le 3n$ always.
\end{proof}

\todo{Proposition~\ref{prop:variance-fpt} was transcribed from version~1 of the arXiv
preprint of \cite{FPT20}; compare with the published version before circulation.}

\begin{remark}\label{rem:variance-fpt}
Only the finiteness of $C_0$ is used, and only in Theorem~\ref{thm:variance}(ii) and
Corollary~\ref{cor:variance-core}. The bound does not depend on the pair costs being
uniform rather than exponential: if the pair costs are independent exponential variables
$E_p$ of rate $1$, then $\defi(p):=1-e^{-E_p}$ is uniform on $[0,1]$ and
$\defi(p)\le E_p$, so every assignment is cheaper for the uniform deficits, and a bound of
the form $c_2n^{1/3}$ for exponential costs implies the same bound for uniform costs.
\end{remark}

For $\eps\in(0,1/2)$ put
\begin{equation}\label{eq:variance-constants}
r_A:=\Bigl\lceil\frac1{2\eps}\Bigr\rceil-1,\qquad
r_B:=\Bigl\lceil\frac1{2\eps}\Bigr\rceil+1,\qquad
\kappa:=\frac{0.3^{\,r_A}}{8\cdot4^{\,r_A+1}},\qquad
C_\eps:=3\Bigl(\Bigl\lceil\frac1{2\eps}\Bigr\rceil+2\Bigr).
\end{equation}
Note that $r_A\ge1$, because $1/(2\eps)>1$. For $K\ge1$ let $q_0(\eps,K)$ be the least
integer $q_0\ge2$ such that every $q\ge q_0$ satisfies
\begin{align}
q^{2\eps}&\ge 4\max\bigl(13(K+5),\,600\bigr)\log q, \tag{Q1}\label{eq:variance-Q1}\\
q^{2\eps}&\ge (20/3)^{\lceil 1/(2\eps)\rceil}, \tag{Q2}\label{eq:variance-Q2}\\
q^{-K}&\ge 3q^3\Bigl[3e^{-0.009(q-1)}+e^{-q^{2\eps}/32}+e^{-\kappa q^{2\eps}}\Bigr]
        +3e^{-0.009q}. \tag{Q3}\label{eq:variance-Q3}
\end{align}
Each of the three conditions holds for all large $q$, so $q_0(\eps,K)$ is finite.

\begin{theorem}\label{thm:variance}
Fix $\eps\in(0,1/2)$ and put $s=q^{-1/2+\eps}$, with $C_\eps$ as in
\eqref{eq:variance-constants}.
\begin{enumerate}
\item[\textup{(i)}] For every $K\ge1$ and every $q\ge q_0(\eps,K)$, with probability at
least $1-150\cdot2^K q^{-K}$ the following two statements hold: no optimal assignment uses
a pair of deficit at least $C_\eps s$; and every team in the support of every optimal
solution of the fractional relaxation has deficit less than $C_\eps s$.
\item[\textup{(ii)}] With $C_0$ as in \eqref{eq:variance-C0}, for every
$q\ge q_0(\eps,2)$,
\begin{align*}
\Var\Tq&\le 4\,C_\eps C_0\,q^{-1/6+\eps}+\frac{7000}{q},\\
\Var\LPq&\le 4\,C_\eps C_0\,q^{-1/6+\eps}+\frac{7000}{q},\\
\Var\Dq&\le 16\,C_\eps C_0\,q^{-1/6+\eps}+\frac{28000}{q}.
\end{align*}
\end{enumerate}
\end{theorem}

The consequence for the core is Corollary~\ref{cor:variance-core} at the end of the
section.

\begin{remark}\label{rem:variance-range}
For $\eps\ge1/2$ part (i) is trivial, since then $C_\eps s\ge 9$ while pair deficits are at
most $1$ and team deficits at most $3$; and the bounds of part (ii) for $\eps\ge1/2$ follow
from the case $\eps=1/4$, because $C_{1/4}\,q^{1/12}=12\,q^{1/12}\le 9\,q^{1/3}\le
C_\eps\,q^{-1/6+\eps}$ for $q\ge4$. The thresholds $q_0(\eps,K)$ are astronomically large
(for $\eps=1/4$, condition \eqref{eq:variance-Q1} alone requires $q>10^9$); the theorem is
asymptotic only.
\end{remark}

\subsection{Preliminaries}\label{sec:variance-prelim}

\subsubsection*{Loads}
Let $\mathcal P$ be the polytope of fractional assignments: the vectors $x\ge0$ indexed by
the $q^3$ teams with $\sum_{e\ni i}x_e=1$ for each agent $i$. Thus
$\LPq=\max_{x\in\mathcal P}\sum_ev(e)x_e$, and the \emph{support} of $x$ is
$\{e:x_e>0\}$. The \emph{load} of a pair $p$ under $x$ is
\[
\ell_p(x):=\sum_{e\supset p}x_e ,
\]
the sum over the teams containing both agents of $p$. Since every team contains three
pairs and $\sum_ex_e=q$,
\begin{equation}\label{eq:variance-loads}
\begin{gathered}
0\le\ell_p(x)\le1,\qquad \sum_p\ell_p(x)=3q,\\
\sum_ev(e)x_e=\sum_p\ell_p(x)\bigl(1-\defi(p)\bigr),\qquad
\sum_p\ell_p(x)\,\defi(p)=3q-\sum_ev(e)x_e .
\end{gathered}
\end{equation}

\subsubsection*{Positions}
A \emph{position} of type $\{X,Y\}$, where $X\ne Y$ are two of the sectors, is a pair
$(x,y)$ with $x\in X$ and $y\in Y$; its \emph{third sector} is the remaining one. Fix
$0<s\le1$. An agent $z$ of the third sector is \emph{$s$-close} to the position $(x,y)$ if
$\defi(x,z)+\defi(y,z)<s$. This event has probability
\[
\pi:=s^2/2\le1/2 ,
\]
the probability that the sum of two independent uniform variables is below $s$. Positions
are \emph{vertex-disjoint} if no two of them share an agent. For a set $S$ of positions of
one type, $N(S)$ denotes the set of agents of the third sector that are $s$-close to at
least one position of $S$. The map $N$ is monotone: $S'\subseteq S$ implies
$N(S')\subseteq N(S)$.

\subsubsection*{Standard facts}
The following are used without proof.

\begin{enumerate}
\item[\textup{(C1)}] \emph{Chernoff bound, lower tail.} If $Y$ is binomial with parameters
$n$ and $\varrho$, and $0\le\delta\le1$, then
$\P(Y\le(1-\delta)n\varrho)\le\exp(-\delta^2n\varrho/2)$; see \cite{BLM13}.
\item[\textup{(C2)}] \emph{Chernoff--Hoeffding bound.} If $Y$ is binomial with parameters
$n$ and $\varrho$, and $0<a<\varrho$, then
$\P(Y\le an)\le\exp\bigl(-n\,\mathrm{KL}(a\,\|\,\varrho)\bigr)$, where
$\mathrm{KL}(a\,\|\,\varrho)=a\log(a/\varrho)+(1-a)\log\bigl((1-a)/(1-\varrho)\bigr)$; see
\cite{BLM13}.
\item[\textup{(C3)}] \emph{Monotonicity.} If $\varrho\ge\varrho_0$ then
$\P(\mathrm{Bin}(n,\varrho)\le t)\le\P(\mathrm{Bin}(n,\varrho_0)\le t)$ for every $t$.
\item[\textup{(M)}] \emph{Bipartite matchings.} If $y\ge0$ is a weighting of the edges of
a bipartite graph in which the total weight of the edges at each vertex is at most $1$,
then the graph has a matching with at least $\sum_{\text{edges}}y$ edges, hence, the size
of a matching being an integer, with at least $\lceil\sum_{\text{edges}}y\rceil$ edges.
(By K\H{o}nig's theorem there is a vertex cover with as many vertices as the largest
matching has edges, and the total weight is at most the number of vertices in any vertex
cover.)
\item[\textup{(LP)}] \emph{Duality.} The fractional relaxation is feasible and bounded, so
its dual,
\[
\min\sum_iu_i\quad\text{subject to}\quad u(e):=u_a+u_b+u_c\ge v(e)\ \text{for every team }
e=(a,b,c),
\]
with $u$ unrestricted in sign, has an optimal solution, and for every optimal $x$ and every
optimal $u$, $x_e>0$ implies $u(e)=v(e)$ (complementary slackness).
\item[\textup{(ES)}] \emph{Efron--Stein inequality.} If $f$ is a square-integrable function
of independent random variables $E_1,\dots,E_m$, then
$\Var f\le\sum_{i=1}^m\E[\Var_if]$, where $\Var_i$ denotes the variance in $E_i$ with the
other coordinates held fixed; see \cite{BLM13}.
\end{enumerate}

\subsection{Uniform expansion}\label{sec:variance-expansion}

\begin{lemma}[Uniform expansion]\label{lem:variance-expansion}
Consider the market with $n\ge2$ agents in each sector and independent uniform pair
deficits. Let $K\ge1$, $0<s\le1$ and $\pi=s^2/2$, and assume
\begin{align}
n\pi&\ge13(K+2)\log n, \tag{U1}\label{eq:variance-U1}\\
n\pi&\ge600\log n. \tag{U2}\label{eq:variance-U2}
\end{align}
Let $\mathcal U(n,s)$ be the event that, for each of the three position types and
\emph{every} nonempty set $S$ of vertex-disjoint positions of that type,
\begin{enumerate}
\item[\textup{(a)}] $|N(S)|\ge0.3\,n\pi\,|S|$ if $|S|\le1/\pi$;
\item[\textup{(b)}] $|N(S)|\ge0.9\,n$ if $|S|\ge3/\pi$.
\end{enumerate}
Then
\[
\P\bigl(\mathcal U(n,s)\text{ fails}\bigr)\le6\,n^{-K}+3\,e^{-0.009\,n}.
\]
\end{lemma}

\begin{proof}
Fix a type, say $\{A,C\}$ with third sector $B$, and a set
\[
S=\{(x_1,y_1),\dots,(x_m,y_m)\}
\]
of $m$ vertex-disjoint positions of this type. An agent
$z\in B$ belongs to $N(S)$ if and only if $\defi(x_i,z)+\defi(y_i,z)<s$ for some $i$.
Because the positions are vertex-disjoint, the $2m$ pairs $(x_i,z)$, $(y_i,z)$ are distinct
for fixed $z$, and pairs with different $z$ are distinct. So the indicators
$\ind\{z\in N(S)\}$, over the $n$ agents $z\in B$, are independent with probability
\[
\pi_m:=1-(1-\pi)^m,
\]
and $|N(S)|$ is binomial with parameters $n$ and $\pi_m$. The number of sets of $m$
vertex-disjoint positions of one type is $\binom nm^2m!\le n^{2m}$.

\emph{Part} (a). Let $1\le m\le1/\pi$. Since $1-\pi\le e^{-\pi}$ and
$t\mapsto(1-e^{-t})/t$ is decreasing on $(0,1]$,
\[
\pi_m\ge1-e^{-\pi m}\ge(1-e^{-1})\,\pi m=:\varrho_0,\qquad 1-e^{-1}=0.63212\ldots
\]
Apply (C3) and then (C1) with $\varrho=\varrho_0$ and with $\delta$ defined by
$(1-\delta)(1-e^{-1})=0.3$, that is, $\delta=0.52540\ldots$:
\begin{align*}
\P\bigl(|N(S)|<0.3\,n\pi m\bigr)&\le\exp\bigl(-\delta^2(1-e^{-1})\,n\pi m/2\bigr)\\
&=\exp(-0.08724\ldots\,n\pi m)\le\exp(-0.079\,n\pi m).
\end{align*}
By \eqref{eq:variance-U1}, $0.079\,n\pi\ge1.027\,(K+2)\log n\ge(K+2)\log n$, so the last
bound is at most $n^{-(K+2)m}$. The union bound over all sets $S$ of this type with
$1\le|S|\le1/\pi$ gives at most
\[
\sum_{m\ge1}n^{2m}\,n^{-(K+2)m}=\sum_{m\ge1}n^{-Km}\le2\,n^{-K},
\]
the last step because $n^K\ge2$.

\emph{Part} (b). Put $m_b:=\lceil3/\pi\rceil$. A set with $|S|\ge3/\pi$ has at least $m_b$
elements, so by monotonicity of $N$ it is enough to treat the sets with $|S|=m_b$ (if
there are none, (b) is vacuous). For such a set,
$\pi_{m_b}\ge1-e^{-\pi m_b}\ge1-e^{-3}=0.95021\ldots$, and by (C3) and (C2),
\begin{align*}
\P\bigl(|N(S)|<0.9\,n\bigr)&\le\exp\bigl(-n\,\mathrm{KL}(0.9\,\|\,1-e^{-3})\bigr)\\
&=\exp(-0.020879\ldots\,n)\le\exp(-0.0208\,n).
\end{align*}
The number of such sets is at most $n^{2m_b}$, and
$2m_b\log n\le(6/\pi+2)\log n$. By \eqref{eq:variance-U2}, $(6/\pi)\log n\le n/100$; and,
using $\pi\le1/2$, $\log n\le n\pi/600\le n/1200$, so $2\log n\le n/600$. The union bound
over the sets of this type with $|S|=m_b$ gives at most
\[
\exp\bigl((0.01+0.001667-0.0208)\,n\bigr)\le\exp(-0.009\,n).
\]

Adding the bounds of parts (a) and (b) and multiplying by $3$ for the three position types
gives the claim.
\end{proof}

The point of Lemma~\ref{lem:variance-expansion} is that the event $\mathcal U(n,s)$
concerns all sets $S$ simultaneously. The set $S$ may therefore be chosen after looking at
the instance (at an optimal assignment, or at an optimal solution of the fractional
relaxation), and no conditioning is involved.

Part (a) of the event stops at $|S|=1/\pi$. The next lemma, which is deterministic, carries
the growth on to linear size.

\begin{lemma}[Growth through the cutoff]\label{lem:variance-growth}
Assume the hypotheses of Lemma~\ref{lem:variance-expansion} and that the event
$\mathcal U(n,s)$ holds, and put $\lambda:=0.3\,n\pi$. Then $\lambda>1$, and for every
nonempty set $S$ of vertex-disjoint positions of one type, with $m=|S|$:
\begin{enumerate}
\item[\textup{(G1)}] $|N(S)|\ge\min(\lambda m,\,0.15\,n)$;
\item[\textup{(G2)}] if $m\ge0.15\,n$, then $|N(S)|\ge0.9\,n$.
\end{enumerate}
Moreover:
\begin{enumerate}
\item[\textup{(G3)}] if real numbers $x_0,x_1,\dots$ satisfy
$x_{i+1}\ge\min(\lambda x_i,\,0.15\,n)$ for all $i\ge0$, then
$x_i\ge\min(\lambda^ix_0,\,0.15\,n)$ for all $i\ge0$.
\end{enumerate}
\end{lemma}

\begin{proof}
By \eqref{eq:variance-U2}, $n\pi\ge600\log2>415$, so $\lambda>1$.

(G1) If $m\le1/\pi$ this is part (a) of $\mathcal U(n,s)$. If $m>1/\pi$, let
$m_0:=\lfloor1/\pi\rfloor$. Then $2\le m_0\le m$, because $\pi\le1/2$, so $S$ has a subset
$S'$ with exactly $m_0$ positions. The set $S'$ consists of vertex-disjoint positions and
$|S'|\le1/\pi$, so by monotonicity and part (a),
\[
|N(S)|\ge|N(S')|\ge0.3\,n\pi\,\lfloor1/\pi\rfloor\ge0.3\,n\pi\,(1/\pi-1)
=0.3\,n\,(1-\pi)\ge0.15\,n,
\]
the last step because $\pi\le1/2$.

(G2) Since $n\pi\ge20$, we have $0.15\,n\ge3/\pi$, so $m\ge3/\pi$ and part (b) of
$\mathcal U(n,s)$ applies.

(G3) By induction on $i$. The case $i=0$ is trivial. If
$x_i\ge\min(\lambda^ix_0,\,0.15\,n)$, then, using $\lambda\ge1$,
\[
x_{i+1}\ge\min(\lambda x_i,\,0.15\,n)
\ge\min(\lambda^{i+1}x_0,\,0.15\,\lambda n,\,0.15\,n)
=\min(\lambda^{i+1}x_0,\,0.15\,n).\qedhere
\]
\end{proof}

There is no rounding issue in Lemmas~\ref{lem:variance-expansion}
and~\ref{lem:variance-growth}: (a), (G1) and (G2) are real lower bounds on integer
cardinalities, and the only integer parts taken are $\lfloor1/\pi\rfloor$ in (G1) and
$\lceil3/\pi\rceil$ in part (b) of Lemma~\ref{lem:variance-expansion}.

\subsection{The support of the fractional relaxation}\label{sec:variance-lpsupport}

\begin{proposition}[Support of the fractional relaxation]\label{prop:variance-lpsupport}
Fix $\eps\in(0,1/2)$ and $K\ge1$, put $s=q^{-1/2+\eps}$ and $\pi=s^2/2=q^{-1+2\eps}/2$, and
let $r_B=\lceil1/(2\eps)\rceil+1$ as in \eqref{eq:variance-constants}. Let $q\ge2$ satisfy
\eqref{eq:variance-Q1} and \eqref{eq:variance-Q2}. Then, on the event $\mathcal U(q,s)$ of
Lemma~\ref{lem:variance-expansion} for the whole market \textup{(}$n=q$\textup{)}, every
optimal solution $x$ of the fractional relaxation and every team $f$ with $x_f>0$ satisfy
\[
\Defi(f)<3\,(r_B+1)\,s .
\]
Consequently this holds with probability at least $1-6\,q^{-K}-3\,e^{-0.009\,q}$.
\end{proposition}

\begin{proof}
By \eqref{eq:variance-Q1}, $n\pi=q^{2\eps}/2\ge\max(13(K+2),600)\log q$, so
Lemma~\ref{lem:variance-expansion} applies with $n=q$, and
Lemma~\ref{lem:variance-growth} holds on $\mathcal U(q,s)$ with
$\lambda=0.15\,q^{2\eps}$. Assume $\mathcal U(q,s)$ from now on.

Fix an optimal $x$, an optimal dual solution $u$ as in (LP), and a support team
$f=(a,b,c)$. Recall that $u(g)\ge v(g)=3-\Defi(g)$ for every team $g$, with equality on
the support of $x$.

\emph{Step 1 (start).} For every $b'\in B$ we have $u(a,b',c)\ge v(a,b',c)$, and
$u(f)=v(f)$. Subtracting,
\[
u_{b'}-u_b\ge\defi(a,b)+\defi(b,c)-\bigl[\defi(a,b')+\defi(b',c)\bigr].
\]
Let $\mathcal B_0:=N(\{(a,c)\})$. For $b'\in\mathcal B_0$ the bracket is below $s$, so
\begin{equation}\label{eq:variance-B1}
u_{b'}>u_b+\defi(a,b)+\defi(b,c)-s\qquad\text{for all }b'\in\mathcal B_0 .
\end{equation}

\emph{Step 2 (one link).} If $g=(a',b',c')$ is a support team and $b''\in B$ is $s$-close
to the position $(a',c')$, then $u(a',b'',c')\ge v(a',b'',c')$ and $u(g)=v(g)$ give
\begin{equation}\label{eq:variance-B2}
u_{b''}-u_{b'}\ge\defi(a',b')+\defi(b',c')-\bigl[\defi(a',b'')+\defi(b'',c')\bigr]>-s,
\end{equation}
because deficits are nonnegative and the bracket is below $s$.

\emph{Step 3 (growth).} Suppose that $\mathcal B_i\subseteq B$ is nonempty. Let $H_i$ be
the bipartite graph with vertex classes $A$ and $C$ whose edges are the positions
$(a',c')$ for which some support team $(a',b',c')$ has $b'\in\mathcal B_i$, and put
\[
y(a',c'):=\sum_{b'\in\mathcal B_i}x_{(a',b',c')} .
\]
Each team has exactly one agent in $B$ and each agent has total weight $1$, so
$\sum_{\text{edges}}y=|\mathcal B_i|$; and for each $a'\in A$,
$\sum_{c'}y(a',c')\le\sum_{e\ni a'}x_e=1$, and likewise at each $c'\in C$. By (M), $H_i$
contains a matching $S_i$ with $|S_i|\ge|\mathcal B_i|$. The set $S_i$ is a set of
vertex-disjoint positions of type $\{A,C\}$, each of which is the $(A,C)$ part of a support
team whose agent in $B$ lies in $\mathcal B_i$. Define
\[
\mathcal B_{i+1}:=N(S_i).
\]
By \eqref{eq:variance-B2}, every $b''\in\mathcal B_{i+1}$ satisfies $u_{b''}>u_{b'}-s$ for
some $b'\in\mathcal B_i$. With \eqref{eq:variance-B1} and induction on $i$,
\begin{equation}\label{eq:variance-B3}
u_{b''}>u_b+\defi(a,b)+\defi(b,c)-(i+1)\,s\qquad\text{for all }b''\in\mathcal B_i
\text{ and all }i\ge0 .
\end{equation}

\emph{Step 4 (sizes).} Part (a) of $\mathcal U(q,s)$ with $|S|=1$ gives
$|\mathcal B_0|\ge\lambda$. By (G1) and $|S_i|\ge|\mathcal B_i|$,
$|\mathcal B_{i+1}|\ge\min(\lambda|\mathcal B_i|,\,0.15\,q)$, so by (G3)
\[
|\mathcal B_i|\ge\min(\lambda^{i+1},\,0.15\,q)\qquad\text{for all }i\ge0;
\]
in particular every $\mathcal B_i$ is nonempty and the construction of Step~3 can be
continued. Put $h:=\lceil1/(2\eps)\rceil=r_B-1$. Then $2\eps(h+1)\ge1+2\eps$, so
\[
\lambda^{h+1}=0.15^{\,h+1}\,q^{2\eps(h+1)}\ge0.15^{\,h+1}\,q^{1+2\eps}\ge0.15\,q,
\]
the last inequality being \eqref{eq:variance-Q2}, that is, $q^{2\eps}\ge(20/3)^h$. Hence
$|\mathcal B_h|\ge0.15\,q$, so $|S_h|\ge0.15\,q$, and (G2) gives
\[
|\mathcal B_{r_B}|=|N(S_h)|\ge0.9\,q .
\]
(The count is: $h=\lceil1/(2\eps)\rceil$ applications of (G1) after $\mathcal B_0$, then
one application of (G2); in total $r_B=h+1$ steps after $\mathcal B_0$.)

\emph{Step 5 (symmetry).} Lemma~\ref{lem:variance-expansion} covers all three position
types, so the same argument with the roles of the sectors permuted gives sets
$\mathcal A_r\subseteq A$, $\mathcal B_r\subseteq B$, $\mathcal C_r\subseteq C$, each of
size at least $0.9\,q$, with $r=r_B$ and
\begin{align*}
u_{a''}&>u_a+\defi(a,b)+\defi(a,c)-(r+1)\,s&&\text{for }a''\in\mathcal A_r,\\
u_{b''}&>u_b+\defi(a,b)+\defi(b,c)-(r+1)\,s&&\text{for }b''\in\mathcal B_r,\\
u_{c''}&>u_c+\defi(a,c)+\defi(b,c)-(r+1)\,s&&\text{for }c''\in\mathcal C_r .
\end{align*}

\emph{Step 6 (close).} The teams whose agent in $A$ lies outside $\mathcal A_r$ carry
total weight $q-|\mathcal A_r|\le0.1\,q$; likewise for $\mathcal B_r$ and $\mathcal C_r$.
Since the total weight of all teams is $q>0.3\,q$, some support team $t=(a'',b'',c'')$
lies in $\mathcal A_r\times\mathcal B_r\times\mathcal C_r$. Adding the three inequalities
of Step~5 and using $u(t)=v(t)\le3$ and $u(f)=v(f)=3-\Defi(f)$,
\[
3\ge v(t)=u(t)>u(f)+2\,\Defi(f)-3(r+1)\,s=3+\Defi(f)-3(r+1)\,s,
\]
so $\Defi(f)<3(r_B+1)\,s$.

The probability statement follows from Lemma~\ref{lem:variance-expansion} with $n=q$.
\end{proof}

The argument uses complementary slackness, so it says nothing about the optimal assignment
when $\Dq>0$. That case needs a different argument.

\subsection{The optimal assignment}\label{sec:variance-assignment}

\begin{proposition}[Pairs used by an optimal assignment]\label{prop:variance-assignment}
Fix $\eps\in(0,1/2)$ and $K\ge1$, put $s=q^{-1/2+\eps}$ and $\pi=s^2/2$, and let $r_A$ and
$\kappa$ be as in \eqref{eq:variance-constants}. Let $q\ge2$ satisfy
\eqref{eq:variance-Q1}. Then the probability that some optimal assignment uses a pair of
deficit at least $(r_A+2)\,s$ is at most
\begin{equation}\label{eq:variance-etaA}
\eta_A:=3q^3\Bigl[6\,(q-1)^{-(K+3)}+3\,e^{-0.009\,(q-1)}+e^{-q^{2\eps}/32}
+e^{-\kappa q^{2\eps}}\Bigr].
\end{equation}
Since $(q-1)^{-(K+3)}\le2^{K+3}q^{-(K+3)}$, the first term of $\eta_A$ is at most
$144\cdot2^K q^{-K}$.
\end{proposition}

\begin{proof}
The law of the market is invariant under permutations of the three sectors, so it is
enough to bound the probability for pairs with one agent in $A$ and one in $B$ and to
multiply by $3$. Fix such a pair $p=(a,b)$, and write $\delta_A:=(r_A+2)\,s$. Note that
\eqref{eq:variance-Q1} and $q\ge2$ force $q^{2\eps}\ge2400\log2$, so $q\ge3$.

\emph{Step 1 (threshold; deterministic).} Let $W_0$ be the least value of
$\Defi(\mathcal M)-\defi(p)$ over the perfect assignments $\mathcal M$ that use $p$, and
$W_1$ the least value of $\Defi(\mathcal M)$ over the perfect assignments that do not use
$p$; both classes are nonempty. Neither $W_0$ nor $W_1$ depends on $\defi(p)$. The least
deficit of a perfect assignment is $\min(W_0+\defi(p),\,W_1)$. If an optimal assignment
uses $p$, then $W_0+\defi(p)\le W_1$, that is, $\defi(p)\le\rho_p:=W_1-W_0$. Hence
\begin{equation}\label{eq:variance-A1}
\{\text{some optimal assignment uses $p$ and }\defi(p)\ge\delta_A\}
\subseteq\{\rho_p\ge\delta_A\}.
\end{equation}

\emph{Step 2 (decoupling; deterministic).} For $c\in C$ let $I_c$ be the market obtained
by removing the agents $a$, $b$ and $c$ (it has $n:=q-1$ agents in each sector), and let
$M_c$ be an optimal assignment of $I_c$, chosen by a fixed rule (say the
lexicographically first optimum), so that $M_c$ is a function of the pair deficits of
$I_c$ only. A perfect assignment that uses $p$ contains the team $(a,b,c)$ for exactly one
$c$, and the rest of it is a perfect assignment of $I_c$. Hence
\begin{equation}\label{eq:variance-A2}
W_0=\min_{c\in C}\bigl[\defi(a,c)+\defi(b,c)+\Defi(M_c)\bigr].
\end{equation}

\emph{Step 3 (rotations; deterministic).} Write $M_c=\{(a_j,b_j,c_j):j=1,\dots,n\}$. For
$k\ge1$ and distinct indices $j_1,\dots,j_k$, the \emph{rotation} of $M_c$ along
$(j_1,\dots,j_k)$ is the following set of $q$ teams: the teams of $M_c$ with index outside
$\{j_1,\dots,j_k\}$; the team $(a_{j_1},b,c_{j_1})$; the teams
$(a_{j_{i+1}},b_{j_i},c_{j_{i+1}})$ for $i=1,\dots,k-1$; and the team $(a,b_{j_k},c)$.
Every agent appears exactly once (in $A$: all $a_j$ and $a$; in $C$: all $c_j$ and $c$; in
$B$: $b$ in the team with index $j_1$, $b_{j_i}$ in the team with index $j_{i+1}$,
$b_{j_k}$ with $a$ and $c$, and the others unmoved), so the rotation is a perfect
assignment, and it does not use $p$, because $b$ is teamed with $a_{j_1}\ne a$. The pairs
it uses are: the pairs $(a_j,c_j)$ for all $j$, and $(a,c)$; the pairs containing an agent
of $B$ in the unmoved teams; and the $2(k+1)$ \emph{new pairs}
\[
(a_{j_1},b),\ (b,c_{j_1});\qquad
(a_{j_{i+1}},b_{j_i}),\ (b_{j_i},c_{j_{i+1}})\quad(1\le i<k);\qquad
(a,b_{j_k}),\ (b_{j_k},c).
\]
Let $\nu$ be the sum of the deficits of the new pairs. The deficit of the rotation is at
most $\defi(a,c)+\Defi(M_c)+\nu\le\defi(a,c)+\defi(b,c)+\Defi(M_c)+\nu$, so
\[
W_1\le\defi(a,c)+\defi(b,c)+\Defi(M_c)+\nu\qquad
\text{for every $c$ and every rotation of $M_c$.}
\]
Taking $c$ to be a minimiser in \eqref{eq:variance-A2} shows that $\rho_p\le\nu$ for every
rotation of that $M_c$. Therefore, with
\[
\mathcal R(p,c):=\{\text{some rotation of $M_c$ has }\nu<\delta_A\},
\]
we have
\begin{equation}\label{eq:variance-A3}
\{\rho_p\ge\delta_A\}\subseteq\bigcup_{c\in C}\mathcal R(p,c)^{\mathrm c}.
\end{equation}
(The minimiser depends on the pairs at $a$ and at $b$; the union over all $q$ fixed values
of $c$ removes that dependence.)

\emph{Step 4 (three independent pieces; $c$ fixed).} On the index set $\{1,\dots,n\}$ of
$M_c$ write $j\to l$ if $b_j$ is $s$-close to the position $(a_l,c_l)$ (the case $j=l$ is
allowed), and define
\begin{align*}
L&:=\{l:\ b_l\text{ is $s$-close to }(a,c)\}=\{l:\ \defi(a,b_l)+\defi(b_l,c)<s\},\\
J&:=\{j:\ b\text{ is $s$-close to }(a_j,c_j)\}=\{j:\ \defi(a_j,b)+\defi(b,c_j)<s\}.
\end{align*}
Three disjoint families of pairs are involved: the pairs of $I_c$, which determine $M_c$
and the relation $\to$; the pairs $(a,b_l)$, $(b_l,c)$, which together with $M_c$ determine
$L$; and the pairs $(a_j,b)$, $(b,c_j)$, which together with $M_c$ determine $J$. Hence,
conditionally on $I_c$, the sets $L$ and $J$ are independent, and each of them contains
each index independently with probability $\pi$. In particular, $J$ is independent of
$(I_c,L)$, with that law.

Let $\Gamma^0:=L$ and $\Gamma^{i+1}:=\Gamma^i\cup\{j:\ j\to l\text{ for some }
l\in\Gamma^i\}$. Thus $\Gamma^i$ is the set of indices from which $L$ is reached by a path
of the relation $\to$ with at most $i$ steps.

\emph{Claim.} If $J\cap\Gamma^{r_A}\ne\emptyset$ then $\mathcal R(p,c)$ holds. Indeed, let
$j$ be in the intersection and take a path $j=j_1\to j_2\to\dots\to j_k\in L$ with the
fewest steps; its vertices are distinct and $k-1\le r_A$. In the rotation along
$(j_1,\dots,j_k)$, the two new pairs at $b$ have deficits summing to less than $s$
(because $j_1\in J$), the two new pairs at $b_{j_i}$ have deficits summing to less than
$s$ for each $i<k$ (because $j_i\to j_{i+1}$), and the two new pairs at $b_{j_k}$ have
deficits summing to less than $s$ (because $j_k\in L$). So
$\nu<(k+1)\,s\le(r_A+2)\,s=\delta_A$.

\emph{Step 5 (expansion on $I_c$).} The market $I_c$ has $n=q-1\ge q/2$ agents in each
sector and independent uniform pair deficits. By \eqref{eq:variance-Q1},
$n\pi\ge q^{2\eps}/4\ge\max(13(K+5),600)\log q$, so Lemma~\ref{lem:variance-expansion}
applies to $I_c$ with $K+3$ in place of $K$:
\[
\P\bigl(\mathcal U(I_c)\text{ fails}\bigr)\le6\,(q-1)^{-(K+3)}+3\,e^{-0.009\,(q-1)},
\]
where $\mathcal U(I_c)$ denotes the event $\mathcal U(n,s)$ for the market $I_c$. On
$\mathcal U(I_c)$, for any nonempty index set $R$ the positions
$\{(a_l,c_l):l\in R\}$ are vertex-disjoint and of type $\{A,C\}$, and, identifying the
agent $b_j$ of $I_c$ with its index $j$,
\[
N\bigl(\{(a_l,c_l):l\in R\}\bigr)=\{j:\ j\to l\text{ for some }l\in R\}.
\]
So by (G1), whenever $L\ne\emptyset$,
\[
|\Gamma^{i+1}|\ge\min(\lambda\,|\Gamma^i|,\,0.15\,n),\qquad\lambda=0.3\,n\pi,
\]
and by (G3), $|\Gamma^{r_A}|\ge\min(\lambda^{r_A}|L|,\,0.15\,n)$. (Only part (a) of the
event and the cutoff step (G1) are used here; part (b) is not needed for
Proposition~\ref{prop:variance-assignment}.)

\emph{Step 6 (probability for fixed $p$ and $c$).} The cardinality $|L|$ is binomial with
parameters $n$ and $\pi$, so by (C1) with $\delta=1/2$,
\[
\P\bigl(|L|<n\pi/2\bigr)\le e^{-n\pi/8}\le e^{-q^{2\eps}/32}.
\]
Let
\[
\mathcal G:=\mathcal U(I_c)\cap\{|L|\ge n\pi/2\},\qquad G:=\Gamma^{r_A}.
\]
Both $\mathcal G$ and $G$ are functions of $(I_c,L)$. On $\mathcal G$, by Step~5,
\begin{align*}
|G|&\ge\min\Bigl(\frac{0.3^{\,r_A}(n\pi)^{r_A+1}}2,\ 0.15\,n\Bigr),\qquad\text{hence}\\
\pi\,|G|&\ge\min\Bigl(\frac{0.3^{\,r_A}(n\pi)^{r_A+2}}{2n},\ 0.15\,n\pi\Bigr).
\end{align*}
Here $n\pi\ge q^{2\eps}/4$ and
$2\eps(r_A+1)=2\eps\lceil1/(2\eps)\rceil\ge1$, so
$(n\pi)^{r_A+1}\ge q^{2\eps(r_A+1)}/4^{r_A+1}\ge q/4^{r_A+1}$. With $n\le q$ this gives
\[
\frac{0.3^{\,r_A}(n\pi)^{r_A+2}}{2n}
=\frac{0.3^{\,r_A}}2\,(n\pi)\,\frac{(n\pi)^{r_A+1}}n
\ge\frac{0.3^{\,r_A}}2\cdot\frac{q^{2\eps}}4\cdot\frac1{4^{r_A+1}}=\kappa\,q^{2\eps}.
\]
Also $0.15\,n\pi\ge0.0375\,q^{2\eps}\ge\kappa\,q^{2\eps}$, since $\kappa\le0.3/128$ for
$r_A\ge1$. Thus $\pi\,|G|\ge\kappa\,q^{2\eps}$ on $\mathcal G$. Since $J$ is independent of
$(I_c,L)$ and contains each index independently with probability $\pi$,
\[
\P\bigl(J\cap G=\emptyset\bigm|I_c,L\bigr)=(1-\pi)^{|G|}\le e^{-\pi|G|}.
\]
By the Claim, $\mathcal R(p,c)^{\mathrm c}\subseteq\{J\cap G=\emptyset\}$, so
\begin{align}
\P\bigl(\mathcal R(p,c)^{\mathrm c}\bigr)
&\le\P(\mathcal G^{\mathrm c})
+\E\bigl[\ind_{\mathcal G}\,\P\bigl(J\cap G=\emptyset\bigm|I_c,L\bigr)\bigr]\notag\\
&\le\P\bigl(\mathcal U(I_c)\text{ fails}\bigr)+\P\bigl(|L|<n\pi/2\bigr)
+e^{-\kappa q^{2\eps}}\notag\\
&\le6\,(q-1)^{-(K+3)}+3\,e^{-0.009\,(q-1)}+e^{-q^{2\eps}/32}+e^{-\kappa q^{2\eps}}.
\label{eq:variance-A4}
\end{align}

\emph{Step 7 (union bounds).} By \eqref{eq:variance-A1} and \eqref{eq:variance-A3}, for
each of the $q^2$ pairs $p=(a,b)$ with $a\in A$ and $b\in B$,
\[
\P\bigl(\text{some optimal assignment uses $p$ and }\defi(p)\ge\delta_A\bigr)
\le\sum_{c\in C}\P\bigl(\mathcal R(p,c)^{\mathrm c}\bigr),
\]
a sum of $q$ terms, each bounded by \eqref{eq:variance-A4}. Summing over the $q^2$ pairs of
this type and multiplying by $3$ for the three types of pair gives $3q^3$ times the bound
\eqref{eq:variance-A4}, which is $\eta_A$.

Finally, $3q^3\cdot6\,(q-1)^{-(K+3)}\le18\cdot2^{K+3}q^{-K}=144\cdot2^Kq^{-K}$.
\end{proof}

\begin{remark}\label{rem:variance-assignment}
(1) Uniqueness of the optimal assignment is not used in
Proposition~\ref{prop:variance-assignment}. (2) When $1/(2\eps)$ is an integer, the set
$\Gamma^{r_A}$ is only guaranteed to have size of order $q$ with a small constant factor,
not $0.15\,n$; this is enough, as Step~6 shows. (3) Lemma~\ref{lem:variance-expansion} is
applied to the $3q^3$ markets $I_c$, one for each pair $p$ and each third agent $c$; this
is why it is invoked with $K+3$ in place of $K$.
\end{remark}

\subsection{A Poincar\'e inequality for the exponential measure}\label{sec:variance-poincare}

The inequality of the next lemma, with the constant $4$, is classical
\cite{BobkovLedoux97}. The short proof is included to keep the section self-contained and
to make the regularity hypotheses explicit.

\begin{lemma}[Poincar\'e inequality]\label{lem:variance-poincare}
Let $\mu$ be the law of $m$ independent exponential variables of rate $1$ on
$[0,\infty)^m$. Let $f\colon[0,\infty)^m\to\R$ be continuous, bounded, and Lipschitz in each
coordinate separately. Then each partial derivative $\partial_if$ exists $\mu$-almost
everywhere, and
\[
\Var_\mu(f)\le4\,\E_\mu\Bigl[\sum_{i=1}^m(\partial_if)^2\Bigr].
\]
\end{lemma}

\begin{proof}
Consider one dimension first. Let $g\colon[0,\infty)\to\R$ be bounded and Lipschitz. It is
absolutely continuous, and $g'$ exists almost everywhere and is bounded. Put $h:=g-g(0)$.
Since $h^2$ is absolutely continuous and $h(0)=0$, integration by parts on $[0,R]$ gives
\[
\int_0^Rh(x)^2e^{-x}\dd x=-h(R)^2e^{-R}+2\int_0^Rh(x)h'(x)e^{-x}\dd x .
\]
Let $R\to\infty$. Since $h$ is bounded the boundary term vanishes, and the Cauchy--Schwarz
inequality gives
\[
\E h^2=2\,\E[hh']\le2\,(\E h^2)^{1/2}(\E h'^2)^{1/2},
\]
that is, $\E h^2\le4\,\E h'^2$, expectations being with respect to the exponential law of
rate $1$. Hence
\[
\Var g\le\E\bigl(g-g(0)\bigr)^2\le4\,\E g'^2 .
\]

Now let $f$ be as in the statement. For every value of the other coordinates, the function
$t\mapsto f(\dots,t,\dots)$ of the $i$-th coordinate is bounded and Lipschitz, so by the
one-dimensional case $\Var_if\le4\,\E_i(\partial_if)^2$, where $\E_i$ is the expectation
in the $i$-th coordinate alone and the derivative exists for almost every $t$. The set
where $\partial_if$ fails to exist is measurable and meets every line parallel to the
$i$-th axis in a null set, so it is $\mu$-null by Fubini's theorem. Taking expectations,
summing over $i$ and applying (ES) gives the claim.
\end{proof}

\subsection{Gradients}\label{sec:variance-gradient}

Write each pair deficit as
\[
\defi(p)=e^{-E_p},
\]
with $(E_p)$ independent exponential variables of rate $1$ indexed by the $m:=3q^2$ pairs.
Then the deficits are independent and uniform on $[0,1]$, as required, and the value of the
pair $p$ is $1-e^{-E_p}$. Regard $\Tq$, $\LPq$ and $\Dq$ as functions of
$E=(E_p)\in[0,\infty)^m$, and write $\nabla$ for the gradient in these coordinates. For a
perfect assignment $\mathcal M$ we write $p\in\mathcal M$ when $\mathcal M$ uses the pair
$p$.

\begin{lemma}[Gradients and uniqueness of the loads]\label{lem:variance-gradient}
\begin{enumerate}
\item[\textup{(1)}] The functions $\Tq$, $\LPq$ and $\Dq$ are continuous and bounded by
$3q$. Each is Lipschitz in each coordinate, with constant $1$ for $\Tq$ and $\LPq$ and
constant $2$ for $\Dq$.
\item[\textup{(2)}] Almost surely the optimal assignment $\Mopt$ is unique, and for every
pair $p$,
\[
\partial\Tq/\partial E_p=\ind\{p\in\Mopt\}\,\defi(p).
\]
\item[\textup{(3)}] Optimal solutions $x$ of the fractional relaxation need not be unique,
but almost surely all of them have the same load vector $\ell^*=(\ell^*_p)$, and for every
pair $p$,
\[
\partial\LPq/\partial E_p=\ell^*_p\,\defi(p).
\]
\item[\textup{(4)}] Consequently, almost surely,
\begin{align*}
|\nabla\Tq|^2&=\sum_{p\in\Mopt}\defi(p)^2\le3q,\\
|\nabla\LPq|^2&=\sum_p(\ell^*_p)^2\,\defi(p)^2\le\sum_p\ell^*_p=3q,\\
|\nabla\Dq|^2&\le2\,|\nabla\Tq|^2+2\,|\nabla\LPq|^2 .
\end{align*}
\end{enumerate}
\end{lemma}

\begin{proof}
For a point $\ell\in\R^m$ write
\[
\varphi_\ell(E):=\sum_p\ell_p\bigl(1-e^{-E_p}\bigr);
\]
it is smooth, with $\partial\varphi_\ell/\partial E_p=\ell_pe^{-E_p}=\ell_p\,\defi(p)$. The
function $\Tq$ is the maximum of $\varphi_\ell$ over the finitely many indicator vectors
$\ell$ of the sets of pairs used by perfect assignments. By \eqref{eq:variance-loads},
$\LPq$ is the maximum of $\sum_ev(e)x_e=\varphi_{\ell(x)}(E)$ over $x\in\mathcal P$, that
is, the maximum of $\varphi_\ell$ over the \emph{load polytope}
$\mathcal L:=\{\ell(x):x\in\mathcal P\}$, a polytope in $[0,1]^m$ (the linear image of
$\mathcal P$). Since $\ell\mapsto\varphi_\ell(E)$ is linear, $\LPq$ is the maximum of
$\varphi_\ell$ over the finitely many vertices of $\mathcal L$.

(1) If $0\le\ell_p\le1$ for all $p$ and $\sum_p\ell_p=3q$, then $\varphi_\ell$ is
$1$-Lipschitz in each coordinate and takes values in $[0,3q]$. A maximum of finitely many
such functions is continuous and has the same two properties; and $\Dq$ is the difference
of two of them, with $0\le\Dq\le3q$.

(2), (3) Let $V$ be a finite set of points of $\R^m$ (the indicator vectors of the
assignments, or the vertices of $\mathcal L$). In the coordinates $X_p:=1-\defi(p)$ we
have $\varphi_\ell=\langle X,\ell\rangle$. If the maximiser of $\varphi_\ell$ over $V$ is
not unique at $X$, then $\langle X,\ell-\ell'\rangle=0$ for two distinct
$\ell,\ell'\in V$; such $X$ form a finite union of hyperplanes, which is Lebesgue-null.
The map $E\mapsto X$ is a diffeomorphism of $(0,\infty)^m$ onto $(0,1)^m$, so the
exceptional set is null in the coordinates $E$ as well, hence of probability zero. Where
the maximiser $\ell^*$ over $V$ is unique, it remains the unique maximiser on a
neighbourhood (finitely many strict inequalities), so the maximum equals
$\varphi_{\ell^*}$ there, and its partial derivatives are $\ell^*_p\,\defi(p)$. For $\Tq$,
the vector $\ell^*$ is the indicator of the pairs used by $\Mopt$, and distinct perfect
assignments use distinct sets of pairs, so $\Mopt$ is unique. For $\LPq$, uniqueness of the
maximising vertex of $\mathcal L$ means that the set of optimal loads is the single point
$\ell^*$ (the optimal face of $\mathcal L$ is the convex hull of the optimal vertices),
that is, every optimal $x$ has $\ell(x)=\ell^*$. Non-uniqueness of $x$ itself is possible,
because the linear map $x\mapsto\ell(x)$ from $\R^{q^3}$ to $\R^{3q^2}$ is not injective;
it is harmless, since the objective and the gradient depend on $x$ only through $\ell(x)$.

(4) follows from (2) and (3), from $\defi(p)\le1$, $0\le\ell^*_p\le1$ and
$\sum_p\ell^*_p=3q$, and from $(u-w)^2\le2u^2+2w^2$.
\end{proof}

\begin{proposition}[Variance bound with the bad events]\label{prop:variance-badevent}
Let $\delta_A,\delta_B>0$. Let $\mathcal G_A$ be the event that no optimal assignment uses
a pair of deficit at least $\delta_A$, and let $\mathcal G_B$ be the event that every team
in the support of every optimal solution of the fractional relaxation has deficit less
than $\delta_B$. Suppose that $\P(\mathcal G_A\text{ fails})\le\eta_A$ and
$\P(\mathcal G_B\text{ fails})\le\eta_B$. Then
\begin{align*}
\Var\Tq&\le4\bigl[\delta_A\,\E\Phiq+3q\,\eta_A\bigr],\\
\Var\LPq&\le4\bigl[\delta_B\,\E\Phiq+3q\,\eta_B\bigr],\\
\Var\Dq&\le8\bigl[(\delta_A+\delta_B)\,\E\Phiq+3q\,(\eta_A+\eta_B)\bigr].
\end{align*}
\end{proposition}

\begin{proof}
By Lemma~\ref{lem:variance-gradient}(1), the functions $\Tq$, $\LPq$ and $\Dq$ satisfy
the hypotheses of Lemma~\ref{lem:variance-poincare}. That lemma, together with parts
(2)--(4) of Lemma~\ref{lem:variance-gradient}, gives
\begin{align*}
\Var\Tq&\le4\,\E\sum_{p\in\Mopt}\defi(p)^2,\qquad
\Var\LPq\le4\,\E\sum_p(\ell^*_p)^2\defi(p)^2,\\
\Var\Dq&\le8\,\E\Bigl[\sum_{p\in\Mopt}\defi(p)^2+\sum_p(\ell^*_p)^2\defi(p)^2\Bigr].
\end{align*}

On $\mathcal G_A$, every pair used by $\Mopt$ has deficit less than $\delta_A$, so
\[
\sum_{p\in\Mopt}\defi(p)^2\le\delta_A\sum_{p\in\Mopt}\defi(p)=\delta_A\,\Phiq .
\]
Off $\mathcal G_A$ the sum is at most $3q$, by Lemma~\ref{lem:variance-gradient}(4). So
\[
\E\sum_{p\in\Mopt}\defi(p)^2\le\delta_A\,\E\Phiq+3q\,\eta_A .
\]

On $\mathcal G_B$: if $\ell^*_p>0$, take any optimal $x$; then $\ell_p(x)=\ell^*_p>0$, so
some support team contains the pair $p$, and $\defi(p)$ is at most the deficit of that
team, which is less than $\delta_B$. With $\ell^*_p\le1$ and \eqref{eq:variance-loads},
\[
\sum_p(\ell^*_p)^2\defi(p)^2\le\delta_B\sum_p\ell^*_p\,\defi(p)=\delta_B\,(3q-\LPq)
\le\delta_B\,(3q-\Tq)=\delta_B\,\Phiq,
\]
using $\LPq\ge\Tq$. Off $\mathcal G_B$ the sum is at most $3q$. So
\[
\E\sum_p(\ell^*_p)^2\defi(p)^2\le\delta_B\,\E\Phiq+3q\,\eta_B .
\]
The three bounds follow.
\end{proof}

\subsection{Proof of Theorem~\ref{thm:variance} and the consequence for the core}
\label{sec:variance-proof}

\begin{proof}[Proof of Theorem~\ref{thm:variance}]
Let $\eps\in(0,1/2)$, $s=q^{-1/2+\eps}$ and $q\ge q_0(\eps,K)$.

(i) Proposition~\ref{prop:variance-assignment}, whose hypothesis is
\eqref{eq:variance-Q1}, gives: off an event of probability at most $\eta_A$, no optimal
assignment uses a pair of deficit at least
\[
\delta_A:=(r_A+2)\,s=\bigl(\lceil1/(2\eps)\rceil+1\bigr)\,s\le C_\eps\,s .
\]
Proposition~\ref{prop:variance-lpsupport}, whose hypotheses are \eqref{eq:variance-Q1} and
\eqref{eq:variance-Q2}, gives: off an event of probability at most
$\eta_B:=6\,q^{-K}+3\,e^{-0.009\,q}$, every support team of every optimal solution of the
fractional relaxation has deficit less than
\[
\delta_B:=3\,(r_B+1)\,s=3\bigl(\lceil1/(2\eps)\rceil+2\bigr)\,s=C_\eps\,s .
\]
The total failure probability is at most
\[
\eta_A+\eta_B\le(144\cdot2^K+6)\,q^{-K}
+3q^3\Bigl[3e^{-0.009(q-1)}+e^{-q^{2\eps}/32}+e^{-\kappa q^{2\eps}}\Bigr]+3e^{-0.009q}.
\]
By \eqref{eq:variance-Q3} the exponentially small terms are together at most $q^{-K}$, so
\[
\eta_A+\eta_B\le(144\cdot2^K+7)\,q^{-K}\le150\cdot2^K q^{-K}.
\]

(ii) Take $K=2$ and apply Proposition~\ref{prop:variance-badevent} with the events and the
values $\delta_A$, $\delta_B$, $\eta_A$, $\eta_B$ of part (i). By
Proposition~\ref{prop:variance-fpt}, $\E\Phiq\le C_0\,q^{1/3}$, so for
$\delta\in\{\delta_A,\delta_B\}$,
\[
\delta\,\E\Phiq\le C_\eps\,q^{-1/2+\eps}\cdot C_0\,q^{1/3}=C_\eps C_0\,q^{-1/6+\eps}.
\]
By \eqref{eq:variance-Q3} with $K=2$, the exponentially small part of $\eta_A$ and that of
$\eta_B$ are each at most $q^{-2}$, so
\[
3q\,\eta_A\le3q\,(576+1)\,q^{-2}=\frac{1731}q,\qquad
3q\,\eta_B\le3q\cdot7\,q^{-2}=\frac{21}q .
\]
Hence
\begin{align*}
\Var\Tq&\le4\,C_\eps C_0\,q^{-1/6+\eps}+\frac{6924}q,\\
\Var\LPq&\le4\,C_\eps C_0\,q^{-1/6+\eps}+\frac{84}q,\\
\Var\Dq&\le16\,C_\eps C_0\,q^{-1/6+\eps}+\frac{14016}q,
\end{align*}
which are within the stated bounds.
\end{proof}

Recall from Section~\ref{sec:model} that the core is nonempty exactly when $\Dq=0$.

\begin{corollary}\label{cor:variance-core}
Let $\eps\in(0,1/2)$. If $\E\Dq\ge q^{-1/12+\eps}$ for all large $q$, then for all large
$q$
\[
\P(\text{the core is nonempty})=\P(\Dq=0)\le\frac{\Var\Dq}{(\E\Dq)^2}
\le16\,C_\eps C_0\,q^{-\eps}+28000\,q^{-5/6-2\eps},
\]
which tends to $0$. In particular, $\P(\Dq=0)\to0$ whenever $\liminf_q\E\Dq>0$.
\end{corollary}

\begin{proof}
Since $\Dq\ge0$ and $\E\Dq>0$, we have $\{\Dq=0\}\subseteq\{|\Dq-\E\Dq|\ge\E\Dq\}$, and
Chebyshev's inequality gives $\P(\Dq=0)\le\Var\Dq/(\E\Dq)^2$. If
$\E\Dq\ge q^{-1/12+\eps}$, then $(\E\Dq)^2\ge q^{-1/6+2\eps}$, and
Theorem~\ref{thm:variance}(ii) with the same $\eps$ gives, for $q\ge q_0(\eps,2)$,
\[
\P(\Dq=0)\le16\,C_\eps C_0\,q^{-\eps}+28000\,q^{-1}\,q^{1/6-2\eps}
=16\,C_\eps C_0\,q^{-\eps}+28000\,q^{-5/6-2\eps}.
\]
If $\liminf_q\E\Dq>0$, then $\E\Dq\ge q^{-1/12+1/24}$ for all large $q$, and the first part
applies with $\eps=1/24$.
\end{proof}

\begin{remark}\label{rem:variance-vacuous}
The restriction to $\eps<1/2$ in Corollary~\ref{cor:variance-core} loses nothing: since
$\E\Dq\le\E\Phiq\le C_0\,q^{1/3}$, the hypothesis $\E\Dq\ge q^{-1/12+\eps}$ cannot hold for
all large $q$ when $\eps>5/12$, so for $\eps\ge1/2$ it is vacuous.
\end{remark}
